\PuzzleChapterIntro{The Violet Moon Rite}{Logic \textbullet\ Truth-Tables \textbullet\ Constraint Satisfaction}{This puzzle applies a strict ``One Lie Per Witness'' constraint. Instead of evaluating every possible world, the optimal strategy involves identifying internal contradictions within a single suspect's testimony to collapse the possibilities.}
\Lore{

On the Violet Moon, the monastery runs its rites by lanternlight and habit.

Four acolytes---Athena, Quill, Kestrel, Sable---were the only ones assigned to the inner corridor that night.
At dawn, the reliquary was empty and the ash on the threshold showed too many careful steps to be an accident.

The prior refuses confessions taken in panic. Instead, the prior uses an old rule of discipline:

\begin{quote}
\textit{Each acolyte will make four claims. Exactly one claim from each is a lie.}
\end{quote}

It is not mercy. It is arithmetic.

}

\Problem

A relic called the \emph{Violet Moon} is stolen.

\textbf{Rules.}

\begin{itemize}
  \item Exactly one of the four suspects stole the Violet Moon.
  \item Each suspect makes \textbf{four} statements.
  \item For \textbf{each} suspect, \textbf{exactly three} of their statements are true and \textbf{exactly one} is false.
\end{itemize}

Whenever a suspect says ``\emph{exactly one of the following two claims is true},'' interpret it literally: one is true and the other is false.

\medskip

Four suspects---\textbf{Athena}, \textbf{Quill}, \textbf{Kestrel}, \textbf{Sable}---give the following statements. (In each line, ``\emph{X did it}'' means ``\emph{X stole the Violet Moon}.'')

\medskip

\textbf{Athena:}

\begin{enumerate}[label=(\arabic*)]
\item Kestrel did it \emph{if and only if} Sable did \emph{not} do it.
\item Quill did it \emph{if and only if} Sable did \emph{not} do it.
\item Sable did it.
\item Exactly one of these is true: ``Athena did it'' and ``Quill did it.''
\end{enumerate}

\medskip

\textbf{Quill:}

\begin{enumerate}[label=(\arabic*)]
\item Athena did it \emph{if and only if} Kestrel did it.
\item Kestrel did it \emph{if and only if} Sable did it.
\item Exactly one of these is true: ``Kestrel did it'' and ``Sable did it.''
\item If Kestrel did it, then Quill did it.
\end{enumerate}

\medskip

\textbf{Kestrel:}

\begin{enumerate}[label=(\arabic*)]
\item If Quill did it, then Kestrel did \emph{not} do it.
\item Exactly one of these is true: ``Quill did it'' and ``Sable did it.''
\item Quill did it \emph{if and only if} Kestrel did \emph{not} do it.
\item If Kestrel did it, then Sable did \emph{not} do it.
\end{enumerate}

\medskip

\textbf{Sable:}

\begin{enumerate}[label=(\arabic*)]
\item If Quill did it, then Sable did \emph{not} do it.
\item If Sable did it, then Quill did it.
\item Exactly one of these is true: ``Athena did it'' and ``Sable did it.''
\item Quill did it \emph{if and only if} Kestrel did it.
\end{enumerate}

\medskip

\VaultDirective{Who stole the Violet Moon?}

\SolutionPage

Let \(A,Q,K,S\) denote the propositions ``Athena did it,'' ``Quill did it,'' ``Kestrel did it,'' ``Sable did it.''
Exactly one of \(A,Q,K,S\) is true.

\textbf{Key pressure point (Quill).}

Quill(2) says \(K \leftrightarrow S\).
Quill(3) says exactly one of \(K\) and \(S\) is true.
Those two claims cannot both be true, so at least one of Quill(2), Quill(3) is false.

But Quill has \emph{exactly one} false statement, so \emph{exactly one} of Quill(2), Quill(3) is false, and Quill(1) and Quill(4) are true.

\textbf{Case 1: Quill(2) is false and Quill(3) is true.}

Then exactly one of \(K,S\) is true.
Since Quill(4) is true, \(K \Rightarrow Q\).
If \(K\) were true, then \(Q\) would also be true, contradicting that there is exactly one thief.
Hence \(K\) is false, so \(S\) is true.

So in this case, \(S\) (Sable) stole the Violet Moon.

\textbf{Case 2: Quill(2) is true and Quill(3) is false.}

Then \(K \leftrightarrow S\) is true, so \(K\) and \(S\) have the same truth value.
They cannot both be true (only one thief), so \(K\) and \(S\) are both false.

Now Quill(1) is true: \(A \leftrightarrow K\).
With \(K\) false, this forces \(A\) false.
Thus the only remaining candidate for the (unique) thief would be \(Q\) true.

But if \(Q\) is true while \(K\) and \(S\) are false, then Athena(1) says
\(K \leftrightarrow \neg S\), i.e.\ \(\text{false} \leftrightarrow \text{true}\), which is false;
Athena(3) says \(S\), which is also false.
So Athena would have at least two false statements, contradicting the Echo-Stamp rule.

Therefore Case 2 is impossible.

\textbf{Conclusion.}

Only Case 1 remains, so \(S\) is true: \textbf{Sable stole the Violet Moon.}

\FinalAnswer{Sable}

\NextPage
